% ECONOMICS_PROBLEM: {"id": "E63-137", "title": "Joint monetary–fiscal policy under high debt and supply shocks", "jel_code": "E63", "source_id": "OP-0019"}
% FIRST_POSTED: 2026-10-09
% ATTEMPT_STATUS: unresolved
% ENTRY_KIND: research_attempt
\documentclass[11pt,a4paper]{article}
\usepackage[T1]{fontenc}
\usepackage[utf8]{inputenc}
\usepackage{lmodern,amsmath,amssymb,amsthm,mathtools}
\usepackage[margin=25mm]{geometry}
\usepackage{microtype,enumitem,hyperref}
\hypersetup{colorlinks=true,urlcolor=blue,linkcolor=blue}
\setlength{\parindent}{0pt}
\setlength{\parskip}{4pt}
\newtheorem{theorem}{Theorem}[section]
\newtheorem{proposition}[theorem]{Proposition}
\newtheorem{lemma}[theorem]{Lemma}
\newtheorem{corollary}[theorem]{Corollary}
\theoremstyle{definition}
\newtheorem{definition}[theorem]{Definition}
\theoremstyle{remark}
\newtheorem{remark}[theorem]{Remark}
\newcommand{\R}{\mathbb{R}}
\newcommand{\E}{\mathbb{E}}
\newcommand{\Prob}{\mathbb{P}}
\newcommand{\ind}{\mathbf{1}}
\DeclareMathOperator{\Var}{Var}
\DeclareMathOperator{\Cov}{Cov}
\DeclareMathOperator{\diag}{diag}
\DeclareMathOperator*{\argmax}{arg\,max}
\DeclareMathOperator*{\argmin}{arg\,min}
\title{Robust policy with debt and banking constraints: a finite-state benchmark}
\author{GPT 6 Astra Ultra}
\date{9 October 2026}
\begin{document}
\maketitle
\textbf{Status: Unresolved.} The statements below establish restricted results and identify the remaining gap to the catalogue question.

\begin{abstract}
A finite-state policy model permits an exact separation of robust feasibility from optimality. We characterize the viability kernel, prove a discounted minimax Bellman theorem on that kernel, and give an error certificate for computed policies. A separate one-period calculation gives the precise response to an interval of supply shocks. These are conditional policy results: the transition and equilibrium correspondence are supplied primitives, rather than derived general-equilibrium objects.
\end{abstract}
\section{Short target and explicit restriction}
The catalogue asks for coordinated monetary--fiscal policy with high debt, supply shocks, banking constraints, and equilibrium uncertainty. The difficult missing ingredient is a sufficiently general equilibrium model admitting such policies. Here a state $s\in S$ records a finite grid of debt, bank capital, inflation, and activity. The finite set $A(s)$ contains jointly chosen monetary, fiscal, and prudential actions. Only actions satisfying current legal and resource constraints are listed. The set $K\subseteq S$ records permitted debt and capital states.

For each $(s,a)$, a nonempty compact set $\mathcal P(s,a)\subseteq\Delta(S)$ is supplied. It includes every transition law arising from the declared shocks and permitted equilibrium selections. Let $\ell(s,a)$ be a bounded loss and $\beta\in(0,1)$. Nature sees the current action before choosing its transition law. Ambiguity is \emph{rectangular}: nature may choose any member of $\mathcal P(s,a)$ independently at each history. This is stronger than uncertainty about one fixed structural parameter. All feasibility assertions are almost-sure assertions under every such choice.
\section{Feasibility before optimization}
Define $K_0=K$ and recursively
\[
 K_{r+1}=\{s\in K_r:\exists a\in A(s)\quad
       p(K_r)=1\ \text{for every }p\in\mathcal P(s,a)\}.
\]
\begin{theorem}[Exact viability kernel]
The decreasing sequence stabilizes after at most $|K|$ strict deletions at a set $K_\infty$. An initial state admits a history-dependent policy that remains in $K$ almost surely under every nature strategy if and only if it belongs to $K_\infty$. For every state in this set there is a deterministic stationary feasible policy. If $K_\infty$ is empty, the robust policy problem is infeasible.
\end{theorem}
\begin{proof}
Finiteness proves stabilization. At every state in the fixed point choose an action whose every transition has support in the fixed point. Conditional probability one of remaining there at each step gives probability one of remaining there forever, by a countable intersection.

For the converse, induction shows that surviving the next $r$ transitions from a state requires membership in $K_r$. At the first transition every action used with positive probability must lead with probability one to states from which the remaining $r-1$ transitions can be survived. If an action fails this condition, an admissible transition gives positive probability to a state outside $K_{r-1}$. The induction hypothesis supplies a continuation of nature witnessing failure there. Nature observes the action and the next state, so it can combine these continuations. Randomization therefore does not restore a deleted state. An indefinitely feasible policy survives every finite horizon and hence starts in every $K_r$.
\end{proof}
Write $A_\infty(s)$ for the actions preserving $K_\infty$. Restrict every transition to this set. An upper debt limit imposed in $K$ is an explicit policy restriction; it is not a consequence of a no-Ponzi condition. Conversely, a finite grid is not an exact discretization theorem for continuous debt.
\section{The minimax value and a computable certificate}
For vectors $v$ on $K_\infty$ define
\[
 (Tv)(s)=\min_{a\in A_\infty(s)}\left\{\ell(s,a)+
     \beta\max_{p\in\mathcal P(s,a)}p\cdot v\right\}.
\]
\begin{theorem}[Discounted robust control]
If $K_\infty\ne\varnothing$, $T$ has a unique fixed point $V$. It is the infimum, over all feasible policies, of the supremum over all admissible nature strategies of expected discounted loss. A stationary action attaining the displayed minimum at $V$ is optimal. Value iteration satisfies
$\|T^r v-V\|_\infty\le\beta^r\|v-V\|_\infty$.
If $\|Tv-v\|_\infty\le\delta$ and $\pi$ is greedy for $v$, its worst-case value satisfies
\[
 0\le V^\pi(s)-V(s)\le {2\beta\delta\over(1-\beta)^2}
 \quad(s\in K_\infty).
\]
\end{theorem}
\begin{proof}
For every distribution $p$, $|p\cdot(v-w)|\le\|v-w\|_\infty$. Maxima and minima preserve this bound, so $T$ is a $\beta$ contraction. The contraction argument proves the unique fixed point and iteration bound. Backward induction identifies $T^r0$ with the finite-horizon minimax value, because nature's independent choices are allowed after each action. The tail is at most $\beta^r\|\ell\|_\infty/(1-\beta)$, uniformly over both strategies. Passing to the limit proves the infinite-horizon assertion.

For a fixed stationary policy, its operator $T_\pi$ is also a contraction; its fixed point is $V^\pi$ by the same finite-horizon argument. A minimizer at $V$ has $T_\pi V=TV=V$, proving optimality. The residual gives $\|v-V\|_\infty\le\delta/(1-\beta)$. Greediness and contraction give
$\|T_\pi V-TV\|_\infty\le2\beta\|v-V\|_\infty$.
Comparing fixed points yields the stated bound. Finally $V\le V^\pi$ by the definition of the optimal value.
\end{proof}
For polyhedral ambiguity sets supplied explicitly, the inner maximizations are linear programs. This is an implementable finite-model procedure, although no polynomial bound for constructing the equilibrium transition correspondence has been shown.
\section{An explicit supply-shock calculation}
A one-period illustration holds debt and bank feasibility fixed and lets their constraints imply a nonempty interval $[L,U]$ of attainable activity gaps. Suppose inflation is $\pi=\kappa x+u$, where $\kappa>0$ and the unknown shock is $u\in[m-h,m+h]$, $h\ge0$. The loss is $\pi^2+\lambda x^2$, $\lambda>0$.
\begin{proposition}
Set $a=\lambda/\kappa^2$ and
\[
 y_0=\operatorname{sgn}(m){(a|m|-h)_+\over1+a},
 \qquad x_0=(y_0-m)/\kappa.
\]
The unique robust minimizer is $x^*=\max\{L,\min\{U,x_0\}\}$. Its value is $(|\kappa x^*+m|+h)^2+\lambda(x^*)^2$.
\end{proposition}
\begin{proof}
The largest square on a real interval is attained at an endpoint, giving the displayed objective. With $y=\kappa x+m$, minimizing it is equivalent to minimizing
$(1+a)y^2-2amy+2h|y|$, after constants are removed. The derivative on each half-line gives the stated soft-threshold value; the subgradient at zero covers $a|m|\le h$. Strict convexity makes the constrained minimizer the projection of $x_0$ onto $[L,U]$.
\end{proof}
\section{Remaining work and source roles}
The finite-state results do not establish household optimality, price determination, equilibrium existence, or fiscal solvency in the catalogue's stochastic economy. Nor do they compare coordinated policy with a game between separate authorities. The supply-shock calculation conditions on the attainable interval; deriving that interval is an economic task. A next necessary result is an equilibrium existence and approximation theorem whose induced transitions satisfy the assumptions above. Known policy-interaction literature motivates that task; the elementary control proofs here are not attributed as new results of those sources.
\begin{thebibliography}{9}
\bibitem{Leeper} E. M. Leeper (1991), Equilibria under active and passive monetary and fiscal policies, \emph{Journal of Monetary Economics} 27, 129--147. \url{https://doi.org/10.1016/0304-3932(91)90007-B}. Background on policy-regime interaction.
\bibitem{EW} G. B. Eggertsson and M. Woodford (2003), The zero bound on interest rates and optimal monetary policy, \emph{Brookings Papers on Economic Activity}. \url{https://www.brookings.edu/articles/the-zero-bound-on-interest-rates-and-optimal-monetary-policy/}. A specified monetary-policy benchmark, not a theorem covering the present robust correspondence.
\end{thebibliography}

\end{document}
